\documentclass[aspectratio=169,11pt]{beamer}

\usetheme{Madrid}
\usecolortheme{default}
\usefonttheme{professionalfonts}

\usepackage{booktabs}
\usepackage{array}
\usepackage{hyperref}
\usepackage{xcolor}

\definecolor{raceblue}{RGB}{33,83,156}
\definecolor{racegold}{RGB}{214,143,0}
\setbeamercolor{structure}{fg=raceblue}
\setbeamercolor{title}{fg=white,bg=raceblue}
\setbeamercolor{frametitle}{fg=white,bg=raceblue}
\setbeamercolor{block title}{fg=white,bg=raceblue}
\setbeamercolor{block body}{bg=blue!3}

\title[RACE Tutor AI]{RACE Tutor AI}
\subtitle{Classical ML Reading Comprehension Tutor}
\author[Abdullah Sajid \& Huzaifa Khalid]{
Abdullah Sajid \;--\; 22i-2346\\
Huzaifa Khalid \;--\; 22i-2332
}
\date{May 2026}

\begin{document}

\begin{frame}
    \titlepage
\end{frame}

\begin{frame}{Presentation Roadmap}
    \tableofcontents
\end{frame}

\section{Project Overview}

\begin{frame}{Project Goal}
    \begin{block}{Objective}
        Build a fully running reading-comprehension tutor on the RACE dataset.
    \end{block}

    \begin{itemize}
        \item Load and preprocess RACE-style multiple-choice questions.
        \item Train \textbf{Model A} to verify and rank answer options.
        \item Build \textbf{Model B} to generate distractors and graduated hints.
        \item Provide a Streamlit app for quiz practice, hints, predictions, and analytics.
    \end{itemize}

    \vspace{0.2cm}
    \textbf{Final app:} \texttt{frontend/app.py}\\
    \textbf{Best Model A artifact:} \texttt{models/model\_a/logreg\_candidate\_rich.pkl}
\end{frame}

\begin{frame}{Backend Implementation Map}
    \begin{tabular}{p{0.32\linewidth} p{0.58\linewidth}}
        \toprule
        \textbf{File} & \textbf{Purpose} \\
        \midrule
        \texttt{preprocessing.py} & Load, flatten, clean, split, and build candidate examples. \\
        \texttt{feature\_engineering.py} & Sparse text features, dense lexical features, cosine features, scaling. \\
        \texttt{model\_a.py} & Model builders, training helpers, answer prediction. \\
        \texttt{train\_model\_a.py} & CLI for Model A training/evaluation. \\
        \texttt{model\_b.py} & Extractive distractor and hint generation. \\
        \texttt{inference.py} & End-to-end quiz construction for the frontend. \\
        \texttt{frontend/app.py} & Streamlit user interface. \\
        \bottomrule
    \end{tabular}
\end{frame}

\section{Dataset and EDA}

\begin{frame}{Dataset Explanation}
    \begin{itemize}
        \item Dataset: \textbf{RACE}, a reading-comprehension dataset based on English exams.
        \item Each row contains:
        \begin{itemize}
            \item article/passage
            \item question
            \item four options: A, B, C, D
            \item correct answer label
        \end{itemize}
        \item Local expected format:
    \end{itemize}

    \begin{center}
        \texttt{id, article, question, A, B, C, D, answer}
    \end{center}

    \begin{block}{HuggingFace Flattening}
        The \texttt{ehovy/race} format stores options as a list. The backend flattens
        this into explicit columns \texttt{A}, \texttt{B}, \texttt{C}, and \texttt{D},
        keeping only examples with exactly four options.
    \end{block}
\end{frame}

\begin{frame}{EDA: Dataset Size and Answer Distribution}
    \begin{table}
        \centering
        \small
        \begin{tabular}{lrrrrr}
            \toprule
            \textbf{Split} & \textbf{Rows} & \textbf{A} & \textbf{B} & \textbf{C} & \textbf{D} \\
            \midrule
            Train & 70,292 & 15,317 & 18,181 & 19,112 & 17,682 \\
            Validation & 8,787 & 1,914 & 2,273 & 2,390 & 2,210 \\
            Test & 8,787 & 1,915 & 2,272 & 2,389 & 2,211 \\
            \bottomrule
        \end{tabular}
    \end{table}

    \vspace{0.2cm}
    \begin{itemize}
        \item Distribution is not perfectly uniform.
        \item In all splits, option C is most frequent and option A is least frequent.
        \item Train percentages: A 21.8\%, B 25.9\%, C 27.2\%, D 25.2\%.
    \end{itemize}
\end{frame}

\begin{frame}{EDA: Text Length Statistics}
    \begin{table}
        \centering
        \small
        \begin{tabular}{lrrr}
            \toprule
            \textbf{Split} & \textbf{Article Mean/Median} & \textbf{Question Mean/Median} & \textbf{Option Mean/Median} \\
            \midrule
            Train & 275.2 / 279 & 10.0 / 10 & 5.7 / 5 \\
            Validation & 274.4 / 278 & 10.0 / 10 & 5.7 / 6 \\
            Test & 273.7 / 278 & 10.0 / 10 & 5.7 / 5 \\
            \bottomrule
        \end{tabular}
    \end{table}

    \begin{block}{Interpretation}
        Articles are much longer than questions and options. This supports using
        sparse text features and similarity features rather than raw text length alone.
    \end{block}
\end{frame}

\section{Statistical Analysis}

\begin{frame}{Statistical Analysis Used}
    \begin{enumerate}
        \item \textbf{Descriptive statistics}: row counts, word-length means, medians.
        \item \textbf{Class distribution analysis}: answer label frequencies.
        \item \textbf{Chi-square goodness-of-fit test}: tested whether answer labels are uniformly distributed.
        \item \textbf{Candidate feature comparison}: compared correct vs incorrect candidate overlap statistics.
    \end{enumerate}

    \vspace{0.2cm}
    \begin{table}
        \centering
        \small
        \begin{tabular}{lrr}
            \toprule
            \textbf{Split} & \textbf{Chi-square statistic} & \textbf{p-value} \\
            \midrule
            Train & 446.12 & $2.26 \times 10^{-96}$ \\
            Validation & 56.12 & $3.96 \times 10^{-12}$ \\
            Test & 55.63 & $5.03 \times 10^{-12}$ \\
            \bottomrule
        \end{tabular}
    \end{table}
\end{frame}

\begin{frame}{Interpretation of Statistical Analysis}
    \begin{itemize}
        \item The chi-square results show answer labels are statistically different from a perfectly uniform distribution.
        \item Option A appears less frequently than the other labels.
        \item This explains why the early option-A-only binary model was limited.
        \item The final model therefore trains on \textbf{all four candidates} instead of only option A.
    \end{itemize}

    \begin{block}{Correct vs Incorrect Candidate Evidence}
        On the training split:
        \begin{itemize}
            \item Correct option coverage in article: 0.6814
            \item Incorrect option coverage in article: 0.6344
            \item Correct options are slightly longer on average: 6.05 vs 5.76 words
        \end{itemize}
    \end{block}
\end{frame}

\section{Preprocessing}

\begin{frame}{Complete Data Preprocessing Steps}
    \begin{enumerate}
        \item Load CSV splits or HuggingFace RACE.
        \item Flatten nested options into A/B/C/D columns.
        \item Keep only examples with exactly four options.
        \item Lowercase all text.
        \item Remove punctuation and normalize whitespace.
        \item Fill missing text with empty strings.
        \item Drop invalid answer labels.
        \item Drop rows with empty articles or questions.
        \item Expand each question into candidate-level examples.
        \item Fit vectorizer/scaler on train only; reuse on validation/test.
    \end{enumerate}
\end{frame}

\begin{frame}{Feature Engineering Pipeline}
    \begin{columns}
        \begin{column}{0.48\linewidth}
            \textbf{Sparse Features}
            \begin{itemize}
                \item CountVectorizer
                \item Binary bag-of-words
                \item Unigrams + bigrams
                \item Max features: 50,000
                \item Minimum document frequency: 2
            \end{itemize}
        \end{column}
        \begin{column}{0.48\linewidth}
            \textbf{Dense Features}
            \begin{itemize}
                \item Length features
                \item Lexical overlap ratios
                \item TF-IDF cosine similarities
                \item Negation flags
                \item Per-question relative option features
                \item StandardScaler fitted on train only
            \end{itemize}
        \end{column}
    \end{columns}
\end{frame}

\section{Model Selection and Training}

\begin{frame}{Model Selection Process}
    \begin{table}
        \centering
        \small
        \begin{tabular}{p{0.35\linewidth} p{0.48\linewidth}}
            \toprule
            \textbf{Model Tried} & \textbf{Observation} \\
            \midrule
            Option-A logistic model & Simple baseline but ignores B/C/D training signal. \\
            Candidate logistic model & Better aligned with answer selection; trains all four options. \\
            Candidate SVM & Similar but slightly weaker/slower in slice experiments. \\
            Rich candidate logistic model & Best validation and test answer-selection accuracy. \\
            \bottomrule
        \end{tabular}
    \end{table}

    \begin{block}{Final Selection}
        Logistic Regression was selected because it is transparent, stable on sparse text
        features, fast enough for this dataset, and produced the best measured result.
    \end{block}
\end{frame}

\begin{frame}[fragile]{Model Training Process}
    \begin{itemize}
        \item Final training command:
    \end{itemize}

    \scriptsize
    \begin{verbatim}
.venv/bin/python -m backend.train_model_a \
  --train data/raw/train.csv \
  --val data/raw/val.csv \
  --test data/raw/test.csv \
  --model logreg_candidate \
  --class-weight none \
  --ngram-max 2 \
  --c 0.3 \
  --output-name logreg_candidate_rich
    \end{verbatim}
    \normalsize

    \begin{itemize}
        \item Saved model: \texttt{models/model\_a/logreg\_candidate\_rich.pkl}
        \item Training examples: one candidate row per option.
        \item Inference: score A/B/C/D and select highest score.
    \end{itemize}
\end{frame}

\section{Testing and Results}

\begin{frame}{Testing and Evaluation Setup}
    \begin{itemize}
        \item Validation and test sets were held out from training.
        \item Vectorizers and scalers were fitted only on train.
        \item Two evaluation views were used:
        \begin{itemize}
            \item \textbf{Binary candidate metrics}: each candidate is correct/incorrect.
            \item \textbf{Four-option answer selection}: highest-scoring option is compared with the true answer.
        \end{itemize}
        \item The app task uses four-option answer selection, so that is the primary metric.
    \end{itemize}
\end{frame}

\begin{frame}{Model A Result Progression}
    \begin{table}
        \centering
        \small
        \begin{tabular}{lcc}
            \toprule
            \textbf{Model} & \textbf{Validation Accuracy} & \textbf{Test Accuracy} \\
            \midrule
            Option-A binary model & 0.3283 & 0.3214 \\
            Candidate model & 0.3530 & 0.3570 \\
            Rich candidate model & \textbf{0.3809} & \textbf{0.3781} \\
            \bottomrule
        \end{tabular}
    \end{table}

    \begin{block}{Interpretation}
        Moving from option-A-only training to all-candidate training improved ranking.
        Adding bigrams, TF-IDF similarity, negation, and relative option features improved it further.
    \end{block}
\end{frame}

\begin{frame}{Final Performance Metrics}
    \begin{table}
        \centering
        \small
        \begin{tabular}{lrrrr}
            \toprule
            \textbf{Split} & \textbf{Binary Acc.} & \textbf{Macro F1} & \textbf{Answer Acc.} & \textbf{Answer F1} \\
            \midrule
            Validation & 0.7446 & 0.4671 & 0.3809 & 0.3798 \\
            Test & 0.7444 & 0.4660 & 0.3781 & 0.3771 \\
            \bottomrule
        \end{tabular}
    \end{table}

    \begin{itemize}
        \item Binary accuracy is high because most candidate options are incorrect.
        \item Macro F1 reveals the imbalance challenge.
        \item Answer-selection accuracy is the most meaningful app metric.
        \item Test accuracy improved from 32.14\% to 37.81\%.
    \end{itemize}
\end{frame}

\section{Model B and Demo}

\begin{frame}{Model B: Distractors and Hints}
    \begin{columns}
        \begin{column}{0.48\linewidth}
            \textbf{Distractor Generation}
            \begin{itemize}
                \item Extracts passage phrases.
                \item Removes stopwords.
                \item Filters correct-answer duplicates.
                \item Ranks by frequency and question relevance.
                \item Returns three distinct distractors.
            \end{itemize}
        \end{column}
        \begin{column}{0.48\linewidth}
            \textbf{Hint Generation}
            \begin{itemize}
                \item Ranks supporting sentences.
                \item Hint 1: broad focus.
                \item Hint 2: masked evidence.
                \item Hint 3: direct supporting sentence.
            \end{itemize}
        \end{column}
    \end{columns}
\end{frame}

\begin{frame}{Streamlit Demo Flow}
    \begin{enumerate}
        \item Open the app:
        \begin{center}
            \texttt{streamlit run frontend/app.py}
        \end{center}
        \item Load a random RACE sample from \texttt{data/raw/test.csv}.
        \item Answer the question in the Quiz View tab.
        \item Compare your answer with Model A's predicted answer.
        \item Reveal Model B hints one at a time.
        \item View generated distractors.
        \item Check session analytics: attempts, accuracy, latency, and history.
    \end{enumerate}
\end{frame}

\section{Conclusion}

\begin{frame}{Conclusion}
    \begin{itemize}
        \item The project is fully running end to end.
        \item Model A uses a classical ML candidate-ranking pipeline.
        \item Model B provides extractive distractors and graduated hints.
        \item The Streamlit frontend integrates quiz loading, prediction, hints, and analytics.
        \item Final test answer-selection accuracy: \textbf{37.81\%}.
    \end{itemize}

    \begin{block}{Future Improvements}
        Direct ranking objectives, stronger semantic features, transformer embeddings,
        and improved answer-type matching for Model B could further improve performance.
    \end{block}
\end{frame}

\begin{frame}
    \centering
    \Huge Thank You

    \vspace{0.5cm}
    \large Questions?
\end{frame}

\end{document}
